\documentclass[11pt]{article}

\usepackage{series}
\usepackage[T1]{fontenc}
\usepackage[utf8]{inputenc}
\usepackage{lmodern}
\usepackage{microtype}
\usepackage{amsmath,amssymb}
\usepackage{booktabs}
\usepackage{enumitem}

\newcommand{\RR}{\mathbb{R}}
\newcommand{\cH}{\mathcal{H}}
\newcommand{\TT}{\mathrm{TT}}

\title{An SPT-Filtered Euclidean TT Model and Its Conditional Perturbative Properties}
\author{Peter Nero}
\date{Fifth edition\\July 2026}
\hypersetup{
  pdftitle={An SPT-Filtered Euclidean TT Model and Its Conditional Perturbative Properties},
  pdfsubject={Conditional Euclidean filter bounds and their quantum-gravity limits}
}

\begin{document}
\maketitle

\begin{abstract}
This paper isolates the mathematically valid part of a proposed spectral
proper-time (SPT) modification of a transverse-traceless graviton propagator.
It studies a Euclidean model in which selected internal lines carry a
heat-kernel filter. A positive lower endpoint in the filter's proper-time
measure gives Gaussian decay, but that endpoint is an additional source or
schedule assumption: it does not follow from a positive internal spectral
gap. For a fixed graph, the filters suppress every ultraviolet loop
direction exactly when the momenta of the filtered lines have full rank on
the graph's loop space. Equivalently, the subgraph made only from unfiltered
lines must contain no cycle. Thus one filtered graviton line is not enough
in a general multiloop graph. We also prove that a nonzero propagator with a
positive Kallen--Lehmann or Stieltjes representation cannot have permanent
Gaussian decay in the same Euclidean spectral variable. The earlier claims
of automatic reflection positivity, causal support, all-loop finiteness,
BRST consistency, and unitarity are therefore withdrawn. The construction
remains useful as a conditional Euclidean filter or regulator model. It is
not, by itself, a physical ultraviolet completion of quantum gravity or a
selected consequence of Modal Triplet Theory.
\end{abstract}

\section*{Revision note for this edition}
\addcontentsline{toc}{section}{Revision note for this edition}
\begin{description}
\item[Supersedes.] \emph{Modal Triplet Theory: From MTT to a UV-Finite,
Unitary Quantum Gravity}, version~4.
\item[Reason.] Version~4 inferred an external Gaussian filter from an
internal gap, converted a positive Laplace representation into a Stieltjes
representation without a valid inversion theorem, and claimed that one
filtered graviton line made every mixed graph ultraviolet finite. Those
steps are false in the stated generality.
\item[Resolution.] This edition retitles and narrows the work to a Euclidean
TT model. It separates the spectral-gap and proper-time-gap assumptions,
proves the exact loop-rank criterion, supplies an explicit two-loop failure
of the earlier argument, and proves the incompatibility between a nonzero
positive spectral propagator and permanent Gaussian decay.
\item[Retained result.] A proper-time support gap does give a precise
heat-kernel factorization and strong Euclidean graph bounds whenever the
filtered momenta span the full loop space.
\item[Remaining boundary.] No current theorem selects this filter from the
physical q79 MTT geometry or proves reflection positivity, Lorentzian causal
support, gauge consistency, unitarity, or ultraviolet completion for the
filtered theory.
\end{description}

\tableofcontents

\section{Purpose and Orientation}

The original proposal combined three logically different objects:
\begin{enumerate}
\item an internal positive operator with a spectral gap;
\item a Euclidean heat-kernel multiplier acting on external momentum; and
\item a Lorentzian gauge theory intended to describe physical gravitons.
\end{enumerate}
Each object is useful, but none automatically supplies the next. A gap in
an internal finite or compact operator controls internal inversion. A
heat-kernel multiplier controls Euclidean high-frequency behavior. A
physical quantum-gravity model must additionally establish a Lorentzian
state space, gauge identities, causal propagation, and a positive physical
spectrum.

This paper therefore asks a smaller and answerable question:
\begin{quote}
Given an explicitly imposed Euclidean filter, exactly which fixed Feynman
graphs does it control, and which physical conclusions do not follow?
\end{quote}
The answer is useful both positively and negatively. It gives a clean
criterion for filtered calculations and prevents a regulator estimate from
being mistaken for a unitarity theorem.

In plain language, a filter placed on one pipe in a multiloop network does
not slow every independent circulation through that network. It succeeds
only when the filtered pipes intercept the entire cycle space. Likewise, a
positive Euclidean smoothing rule is not automatically a positive spectrum
of physical particles. The two main theorems below make those distinctions
exact.

\section{The Euclidean TT Filter}

\subsection{Model propagator}

Work first in a Euclidean momentum chart. Away from zero momentum, let
\(\Pi_{\TT}(p)\) denote the usual transverse-traceless tensor projector and
consider
\begin{equation}
 \Delta_0(p)
 =\frac{\Pi_{\TT}(p)}{p^2+m_{\rm IR}^2}.
 \label{eq:bare}
\end{equation}
The nonnegative \(m_{\rm IR}\) is only an infrared regulator. It is not a
claim that the physical graviton is massive. Let \(f:[0,\infty)\to\RR\) be a
bounded filter and define
\begin{equation}
 \Delta_f(p)=f(p^2)\Delta_0(p).
 \label{eq:filtered}
\end{equation}
The Gaussian example is \(f(x)=e^{-\tau_0x}\), with \(\tau_0>0\).

Equation~\eqref{eq:filtered} is a model definition. To derive it from MTT,
one would need a selected source map whose compression of the physical
operator equals this multiplier while preserving the required gauge and
causal structures. No such theorem is assumed here.

\subsection{Proper-time support}

\begin{definition}[SPT support-gap hypothesis]
\label{def:spt}
The filter satisfies the SPT support-gap hypothesis if there is a finite
positive measure \(\mu\) and a number \(\tau_0>0\) such that
\begin{equation}
 f(x)=\int_{[\tau_0,\infty)}e^{-tx}\,\mu(dt).
 \label{eq:laplace}
\end{equation}
\end{definition}

\begin{proposition}[Exact SPT factorization]
\label{prop:factor}
Let \(L\geq0\) be self-adjoint and let \(f\) satisfy
Definition~\ref{def:spt}. Then
\begin{equation}
 f(L)=e^{-\tau_0L/2}\,f_0(L)\,e^{-\tau_0L/2},
 \qquad
 f_0(x)=\int_{[\tau_0,\infty)}e^{-(t-\tau_0)x}\,\mu(dt),
 \label{eq:factor}
\end{equation}
where \(f_0(L)\) is bounded and positive. Moreover,
\begin{equation}
 0\leq f(x)\leq \mu([\tau_0,\infty))e^{-\tau_0x}.
 \label{eq:gaussian}
\end{equation}
\end{proposition}

\begin{proof}
The spectral functional calculus permits the factor
\(e^{-\tau_0L}\) to be removed from each integrand in
\eqref{eq:laplace}. Splitting that factor symmetrically gives
\eqref{eq:factor}. Positivity follows from positivity of \(\mu\), and
\eqref{eq:gaussian} follows because \(t\geq\tau_0\).
\end{proof}

The proposition is exact, but its premise matters. A discrete update
schedule can be chosen so that its Laplace measure is supported on
\(\{n\tau_0:n\geq1\}\). That construction proves the implication
\[
 \text{declared schedule}\quad\Longrightarrow\quad
 \text{proper-time support gap}.
\]
It does not prove that the physical MTT source selects that schedule.

\section{An Internal Gap Does Not Produce External Damping}

The earlier paper treated an internal spectral gap and an external
proper-time support gap as though they were the same statement. A simple
compression shows why they are not.

\begin{theorem}[Gap-to-filter implication fails]
\label{thm:gap-fails}
Let \(A_{\rm int}\) be positive self-adjoint on \(\cH_{\rm int}\), and suppose
it has a normalized eigenvector \(e_0\) with eigenvalue
\(\lambda_\ast>0\). For the external spectral variable \(x\geq0\), define
\[
 A(x)=xI+A_{\rm int},
 \qquad
 Bu=\langle e_0,u\rangle.
\]
Then the compressed resolvent is
\begin{equation}
 BA(x)^{-1}B^\ast=\frac{1}{x+\lambda_\ast}.
 \label{eq:compression}
\end{equation}
It has only algebraic decay. Consequently, a positive internal gap alone
does not imply an external factor \(e^{-\tau_0x}\) with \(\tau_0>0\).
\end{theorem}

\begin{proof}
Since \(A(x)e_0=(x+\lambda_\ast)e_0\), compression to the span of \(e_0\)
gives \eqref{eq:compression}. For every \(C,\tau_0>0\),
\[
 \frac{1}{x+\lambda_\ast}
 \not\leq
 C\frac{e^{-\tau_0x}}{x+\lambda_\ast}
\]
for all sufficiently large \(x\). Thus the gap does not supply the claimed
Gaussian factor.
\end{proof}

The internal gap can still be valuable: it controls inverse norms and
decoupling errors. What it cannot do without another theorem is determine
the ultraviolet dependence of the external momentum variable.

\section{Exactly Which Graphs Are Gaussian-Controlled}

\subsection{Loop-momentum map}

Let a connected graph \(\Gamma\) have \(L\) independent loop momenta
\(\ell=(\ell_1,\ldots,\ell_L)\in(\RR^d)^L\). For every filtered internal
line \(j\), write
\begin{equation}
 k_j(\ell,Q)=\sum_{a=1}^L C_{ja}\ell_a+q_j(Q),
 \label{eq:line-momentum}
\end{equation}
where \(Q\) collects external momenta. If line \(j\) has Gaussian parameter
\(\tau_j>0\), the combined damping is
\[
 \exp\!\left(-\sum_j\tau_j|k_j(\ell,Q)|^2\right).
\]
The scalar matrix \(C_\Gamma=(C_{ja})\) records how the filtered lines see the
loop space.

\begin{theorem}[Filtered-loop rank criterion]
\label{thm:rank}
Assume Euclidean propagator denominators are infrared regulated and vertices
grow at most polynomially in loop momenta. The Gaussian factors suppress
every ultraviolet loop direction if and only if
\begin{equation}
 \rank C_\Gamma=L.
 \label{eq:full-rank}
\end{equation}
When \eqref{eq:full-rank} holds, the fixed graph integral is absolutely
ultraviolet convergent. When it fails, the filters alone provide no decay
along at least one nonzero loop direction; ordinary propagator power counting
and renormalization must decide that direction.
\end{theorem}

\begin{proof}
Let \(W=\operatorname{diag}(\tau_j)\). If \(C_\Gamma\) has full column rank,
\(C_\Gamma^{T}WC_\Gamma\) is positive definite. For fixed \(Q\), completing
the square, or using
\(|a+b|^2\geq\tfrac12|a|^2-|b|^2\), gives constants \(c>0\) and \(K_Q<\infty\)
such that
\[
 \sum_j\tau_j|k_j(\ell,Q)|^2
 \geq c|\ell|^2-K_Q.
\]
Hence the integrand is bounded by a polynomial times \(e^{-c|\ell|^2}\),
which is integrable on \(\RR^{dL}\).

If the rank is smaller, choose \(0\neq v\in\ker C_\Gamma\). Along
\(\ell=\ell_0+tv\), every filtered line momentum is independent of \(t\), so
the Gaussian contributes no decay as \(|t|\to\infty\). This does not force
divergence, but it invalidates any convergence proof based only on the
filters.
\end{proof}

\begin{corollary}[Cycle formulation]
\label{cor:cycle}
Let \(D\) be the set of filtered edges and \(U\) the set of unfiltered edges.
For the standard graph momentum routing, \(C_\Gamma\) has full loop rank if
and only if the subgraph spanned by \(U\) is a forest.
\end{corollary}

\begin{proof}
The loop space is the cycle space of \(\Gamma\). The map \(C_\Gamma\) is its
restriction to the filtered edges. Its kernel consists precisely of
circulations supported entirely on \(U\). Such a nonzero circulation exists
exactly when the unfiltered subgraph contains a cycle.
\end{proof}

This formulation makes the earlier error visible. ``At least one filtered
line in the graph'' is much weaker than ``no loop can be routed entirely
through unfiltered lines.''

\subsection{Two worked examples}

\begin{example}[One loop]
For one loop momentum \(k\), a filtered line carrying \(k+Q\) gives
\(C_\Gamma=(1)\). The rank is one, so the Gaussian controls the single
ultraviolet direction.
\end{example}

\begin{example}[Two-loop failure and repair]
Let the independent loop momenta be \(k\) and \(q\). If only a line carrying
\(k\) is filtered, then
\[
 C_\Gamma=\begin{pmatrix}1&0\end{pmatrix},
 \qquad \rank C_\Gamma=1<2.
\]
Along \(k=0\), \(q\to\infty\), the filter is constant. This is an explicit
counterexample to the version~4 claim.

If lines carrying \(k\) and \(k+q\) are both filtered, then
\[
 C_\Gamma=
 \begin{pmatrix}
 1&0\\
 1&1
 \end{pmatrix},
 \qquad \det C_\Gamma=1.
\]
Both loop directions are then Gaussian-controlled.
\end{example}

For mixed gravity--matter diagrams, Corollary~\ref{cor:cycle} says that every
cycle made solely from unfiltered matter lines remains outside the SPT
argument. Such subgraphs require the usual BPHZ, Epstein--Glaser, or EFT
renormalization analysis \cite{Weinberg1965,EpsteinGlaser1973,
BrunettiFredenhagen2000}.

\section{Laplace Positivity Is Not Spectral Positivity}

\subsection{Two different integral representations}

A completely monotone Euclidean function has a positive Laplace
representation,
\begin{equation}
 F(x)=\int_0^\infty e^{-tx}\,\mu(dt).
 \label{eq:laplace-general}
\end{equation}
A Stieltjes two-point function has the more restrictive form
\begin{equation}
 D(x)=\int_0^\infty\frac{\rho(ds)}{x+s},
 \qquad \rho\geq0.
 \label{eq:stieltjes}
\end{equation}
The second form is the Euclidean counterpart of a positive
Kallen--Lehmann spectral representation for a scalar physical component
\cite{Kallen1952,Lehmann1954}. Every Stieltjes function is completely
monotone, but the converse is false \cite{Widder1941,Schilling2012}.

The elementary function \(e^{-x}\) already proves strictness: it is the
Laplace transform of a point mass at \(t=1\), yet it cannot have the
Stieltjes form \eqref{eq:stieltjes} with nonzero positive measure.

\begin{theorem}[Positive spectral tail excludes permanent Gaussian decay]
\label{thm:no-go}
Let
\[
 D(x)=\int_0^\infty\frac{\rho(ds)}{x+s},
 \qquad x>0,
\]
for a nonzero positive measure \(\rho\). Then there do not exist
\(C,\tau,\lambda>0\) such that
\begin{equation}
 D(x)\leq C\frac{e^{-\tau x}}{x+\lambda}
 \label{eq:forbidden}
\end{equation}
for all sufficiently large \(x\).
\end{theorem}

\begin{proof}
Choose \(R<\infty\) for which \(m=\rho([0,R])>0\). Positivity gives
\[
 D(x)\geq\frac{m}{x+R}.
\]
Combining this with \eqref{eq:forbidden} would imply
\[
 \frac{m}{C}
 \leq e^{-\tau x}\frac{x+R}{x+\lambda}\longrightarrow0,
\]
which is impossible.
\end{proof}

\begin{remark}[Scope of the no-go]
The theorem does not prove that every nonlocal quantum theory is nonunitary.
It proves that one specific argument is unavailable: a permanently
Gaussian-damped nonzero propagator cannot obtain positive physical spectral
weight from the same Stieltjes representation. Gauge theories and tensor
fields require additional attention to constraints and physical
cohomology, making it even more important not to infer unitarity from the
scalar Laplace formula alone.
\end{remark}

Osterwalder--Schrader reconstruction requires reflection positivity together
with the other Euclidean axioms \cite{OsterwalderSchrader1973,
OsterwalderSchrader1975}. Positivity of the heat kernel in
\eqref{eq:laplace-general} is not a substitute for that condition.

\section{What the Filter Does Not Prove}

\subsection{Causal support}

Euclidean rapid decay is a statement about large positive covectors. A
Lorentzian retarded kernel requires a hyperbolic operator and a specified
boundary-value prescription. Analytic continuation of an entire Gaussian
multiplier can grow in timelike directions, and nonlocal operators need not
have light-cone-supported retarded kernels. Causal support must therefore be
proved independently; it does not follow from
Proposition~\ref{prop:factor}.

\subsection{BRST and gauge consistency}

The TT projector describes the free physical tensor sector after constraints
have been handled. In an interacting theory, modifying selected propagator
lines after gauge fixing need not preserve the Slavnov--Taylor identities or
the quantum master equation. A legitimate filtered gauge theory would need
one of the following:
\begin{itemize}
\item a gauge-invariant or BV-compatible filtered action;
\item a regulator construction with controlled symmetry breaking and a
proved regulator-removal limit; or
\item an independent physical-state construction in which the filtered
observables act on positive BRST cohomology.
\end{itemize}
Standard BV and perturbative algebraic QFT results remain available for the
local effective theory under their usual hypotheses
\cite{BV1981,BV1983,BarnichBrandtHenneaux2000,HollandsWald2001,
HollandsWald2002,Rejzner2016}. Importing those results does not establish
them for the permanent SPT modification.

\subsection{Matter subgraphs and all-orders language}

If the unfiltered subgraph contains a loop, the SPT filter does not control
that loop. Pure matter subgraphs and rank-deficient mixed subgraphs must be
renormalized by standard local methods. Even when every fixed graph is
absolutely convergent, this does not by itself establish convergence of the
sum over graphs, nonperturbative existence, reflection positivity, or a
unitary Lorentzian scattering operator.

\section{The Infrared Statement That Does Survive}

Suppose \(f\) is differentiable at zero and normalized by \(f(0)=1\). Then
\[
 f(p^2)=1+f'(0)p^2+O(p^4),
\]
so
\begin{equation}
 \Delta_f(p)
 =\Delta_0(p)\left(1+f'(0)p^2+O(p^4)\right).
 \label{eq:ir}
\end{equation}
For \(f(p^2)=e^{-\tau_0p^2}\), the first correction is
\(-\tau_0p^2\). Thus the filter can be chosen to preserve the free
massless-pole residue in a low-momentum expansion.

Equation~\eqref{eq:ir} is only an infrared matching statement for the model
propagator. Recovering General Relativity also requires the nonlinear
vertices, universal conserved stress coupling, gauge identities, and
physical Newton normalization. None follows from \(f(0)=1\).

\section{Placement Inside the Current MTT Program}

The current MTT corpus contains an exact finite internal transverse-traceless
support result on its selected helicity-two carrier. It also contains a
conditional coherent reduction to the four-dimensional Einstein action.
Those are distinct from the SPT filter studied here:
\begin{center}
\begin{tabular}{p{0.30\linewidth}p{0.22\linewidth}p{0.36\linewidth}}
\toprule
Object & Current status & What it does not yet provide\\
\midrule
Finite internal TT support & exact on the declared finite carrier &
Lorentzian graviton pole, Newton normalization, or full stress response\\
Controlled coherent GR reduction & conditional on a higher-dimensional
Einstein sector, truncation, and gap assumptions & derivation of those
inputs from the upper MTT source\\
SPT proper-time endpoint \(\tau_0\) & model assumption or declared schedule &
selection by q79 geometry or identification with a physical Planck scale\\
Filtered graph bounds & exact under Theorem~\ref{thm:rank} &
reflection positivity, BRST consistency, causality, or a sum over graphs\\
\bottomrule
\end{tabular}
\end{center}

In particular, a dimensionless internal eigenvalue such as the selected TT
value recorded in the finite carrier is not automatically a physical mass or
cutoff. A units-and-normalization map is required before it can be compared
with the Planck scale. The broader paper \emph{Perturbative Coherent-Sector
Quantum Gravity and the Heterotic UV-Completion Boundary} should use the
present paper only for the SPT model and its no-go results; its physical
ultraviolet discussion belongs to the separate q79 heterotic worldsheet
contract.

\section{Claim Ledger and Conclusion}

\begin{center}
\begin{tabular}{p{0.47\linewidth}p{0.18\linewidth}p{0.23\linewidth}}
\toprule
Claim & Verdict & Reason\\
\midrule
SPT support gap implies Gaussian Euclidean damping & proved &
Proposition~\ref{prop:factor}\\
Internal spectral gap implies an SPT support gap & false in general &
Theorem~\ref{thm:gap-fails}\\
One filtered graviton line makes every graph finite & withdrawn &
Two-loop rank counterexample\\
Full filtered-loop rank makes a fixed graph UV convergent & proved &
Theorem~\ref{thm:rank}\\
Positive Laplace kernel implies positive spectral density & false in general &
Laplace and Stieltjes classes differ\\
Positive spectral density coexists with permanent Gaussian decay & impossible
for a nonzero propagator & Theorem~\ref{thm:no-go}\\
SPT alone proves BRST, causality, OS positivity, or unitarity & not proved &
Each requires an independent construction\\
SPT is the selected MTT ultraviolet completion & not established &
No same-source selection theorem\\
\bottomrule
\end{tabular}
\end{center}

The corrected conclusion is narrower but firmer. SPT is a useful Euclidean
filter model with an exact support-gap factorization and an exact graph-rank
criterion. It can regularize selected calculations and expose which loop
directions remain uncontrolled. It does not turn an internal MTT gap into a
physical ultraviolet completion. The route from MTT geometry to quantum
gravity must instead supply a selected Lorentzian operator or worldsheet
theory together with its gauge, positivity, causal, normalization, and
nonperturbative contracts.

\bibliographystyle{unsrt}

% BEGIN MTT MANAGED COMPUTATIONAL EVIDENCE
\section*{Computational Evidence and Reproducibility}
The SPT-filtered transverse-traceless model has the perturbative properties proved in this paper only. The open strict-upgrade ledger does not establish reflection positivity, nonperturbative unitarity, or ultraviolet-complete quantum gravity.

The referenced rows are frozen to the curated results repository state identified below. Hashes are grouped in eight-character blocks for line breaking.
\begin{quote}\small
\textbf{Repository:} \url{https://github.com/PeterNero/mtt-results-repro}\\
\textbf{Commit:} \texttt{31247ebb 5c22f3fb b5443024 365433c6 ee0bff4a}\\
\textbf{Manifest:} \path{release/result_manifest.json}\\
\textbf{Manifest SHA-256:}\\
\texttt{fb399689 60b00584 631dbf53 1a708e18}\\
\texttt{ef928d6b 6d935119 c185d7f6 32b1e7cd}
\end{quote}

Tier labels are quoted verbatim from that manifest. A row used directly supports only the specific computational statement identified above; a corpus-state cross-check does not prove this paper's local theorems; and an open row is evidence of an unresolved obligation, never of closure.

\par\begingroup\dimen0=\pagegoal\advance\dimen0 by -\pagetotal\ifdim\dimen0<7\baselineskip\newpage\fi\endgroup
\paragraph{Open boundary (not evidence of closure).}
\begin{itemize}
\setlength{\itemsep}{0.25em}
\item \path{A05/strict_upgrade_ledger} (\emph{open}).\par Current 2/9 strict no-knob upgrade ledger.
\end{itemize}

No imported row changes theorem ownership or promotes a neighboring claim: all local statements retain their stated hypotheses, domains, and limitations.
% END MTT MANAGED COMPUTATIONAL EVIDENCE

\bibliography{main}

\end{document}
